\chapter{CONCLUSIONES}

\section{Conclusiones del trabajo}

Con el experimento ya ejecutado, las conclusiones son:

\begin{itemize}
    \item El pivote del ADR-004 llevó el proyecto a una pregunta que la literatura no tenía respondida, y los datos disponibles bastaron para responderla.
    \item El \textit{pipeline} de preprocesado reconstruido en el paquete \texttt{microraman} reproduce el protocolo de referencia del director con una diferencia máxima absoluta de~0. Las diferencias que se observan entre métricas vienen, por tanto, del fenómeno medido.
    \item Las siete métricas y el motor de \textit{matching} están implementados en dos versiones, pandas/NumPy y PySpark, verificadas entre sí, y se han ejecutado ya sobre el experimento completo con la biblioteca real. La validación contra el protocolo del director alineó las 1885 comparaciones previstas.
    \item El ganador depende del eje de degradación: Pearson con acumulación real (61,5~\% de aciertos), y el grupo coseno, euclídea, HQI y SAM con degradación paramétrica (87,5~\%). Manhattan y Spearman quedan las últimas en los dos ejes. Ninguna diferencia por pares sobrevive la corrección de Holm, así que el orden debe leerse como tendencia y no como superioridad demostrada: solo tres polímeros tienen verdad-terreno disponible.
    \item La comprobación cruzada matiza el \textit{ranking} en lugar de confirmarlo en bloque. Que el orden cambie según el mecanismo de degradación sugiere que la elección de métrica no es universal, sino que depende de cómo se pierda la calidad de la señal.
    \item Dos de los cinco polímeros medidos, BioPBS y ECOVIO, obtienen un 0~\% de acierto con todas las métricas y en todos los niveles, porque la biblioteca de Hagelskjær no los recoge en su catálogo. Es un límite de la base de datos de referencia, no del método desarrollado, y se declara explícitamente en vez de omitirse.
    \item A la escala de esta corrida, 13195 y 8120 filas, distribuir el cómputo con Spark todavía no compensa frente a pandas. El punto de equilibrio exacto queda por medir.
\end{itemize}

\section{Conclusiones personales}

La realización de este TFM me ha dejado aprendizajes técnicos y personales que van más allá de la implementación del \textit{pipeline}. El primero, y quizá el más estructural, es que un proyecto de Big Data no tiene por qué conservar el planteamiento tecnológico con el que se concibió, sobre todo cuando en el desarrollo surgen preguntas nuevas o se dispone de datos que permiten abordar el problema desde otro ángulo. El planteamiento inicial se apoyaba en una arquitectura de ingesta con Apache Kafka, procesamiento en Spark Structured Streaming y clasificación con los modelos preentrenados de Open Specy. Durante el desarrollo decidí, junto con el director, pivotar (ADR-004) hacia el estudio de la robustez de las métricas de similitud frente a la degradación de la calidad espectral. El cambio supuso renunciar a buena parte de la propuesta original, pero permitió centrar el trabajo en una pregunta directamente vinculada a la fiabilidad del \textit{matching} espectral, y abordable de forma experimental con los datos ya disponibles.

Desde el punto de vista técnico, el aprendizaje más claro ha sido comprobar empíricamente que recurrir a una tecnología distribuida no garantiza, por sí solo, un mejor rendimiento. Implementé el motor de \textit{matching} en dos versiones, una de referencia en pandas/NumPy y otra distribuida en PySpark, lo que permitió comparar ambas aproximaciones sobre el mismo problema. A la escala de los experimentos realizados --- 13.195 y 8.120 filas, según el escenario --- el sobrecoste de distribuir el cómputo superó la ganancia obtenida, y pandas resultó más eficiente que Spark. Esto me obligó a distinguir con más precisión entre usar una tecnología Big Data y diseñar una solución que escale de forma eficiente para el volumen real del problema: son cosas distintas, y la primera no implica la segunda.

Otro aprendizaje relevante ha sido dimensionar el peso que tiene la cobertura de los datos de referencia sobre el resultado final. Al ejecutar el experimento observé que BioPBS y ECOVIO obtenían un 0~\% de acierto con las siete métricas, en todos los niveles de calidad. Mi primera hipótesis fue que el problema estaba en el \textit{pipeline} o en las propias métricas; la revisión posterior mostró que ninguno de los dos polímeros figura en la \emph{Hagelskjær Microplastic Solution Library 1.1}, de modo que no existía una referencia correcta contra la que identificarlos. El episodio me enseñó a no interpretar un resultado inesperado como un fallo del algoritmo antes de revisar los datos, las referencias empleadas y las condiciones del experimento: aquí la limitación estaba en la base de datos, no en el método, y acabó condicionando el alcance del análisis estadístico a los tres polímeros con verdad-terreno disponible (PET, PLA, PP).

Trabajar en un problema interdisciplinar ha sido, en sí mismo, otro aprendizaje. El proyecto no se reducía a implementar algoritmos o levantar un \textit{pipeline}: exigía traducir un problema de identificación de microplásticos por espectroscopía Raman a un problema computacional medible y reproducible. Por eso fue crítico respetar el protocolo de preprocesado que me proporcionó el director y comprobar, paso a paso, que mi implementación lo reproducía correctamente; el \textit{pipeline} final iguala ese protocolo de referencia con una diferencia máxima absoluta de~0, lo que permitió comparar las métricas sin que ninguna diferencia en el preprocesado contaminara la comparación. De ahí extraigo una distinción que no tenía tan clara al empezar: la validación técnica del código y la validación científica del procedimiento son necesarias las dos, pero no son lo mismo.

En retrospectiva, habría adelantado el análisis de cobertura de la biblioteca de referencia a una fase mucho más temprana del proyecto, antes de ejecutar la matriz experimental completa: saber de antemano qué polímeros contaban con una referencia válida habría permitido separar, ya en el diseño, los casos evaluables de aquellos en los que la ausencia de referencia hacía imposible la identificación. También habría diseñado desde el principio, y no a partir de una sugerencia del director casi al cierre del experimento, un análisis que no se limitara a comprobar si la referencia correcta queda en primer lugar, sino que midiera además el margen de similitud entre la referencia correcta y las incorrectas. Esa extensión queda planteada como trabajo futuro y es ejecutable directamente sobre los datos ya generados.

El propio resultado experimental me ha enseñado a leer con cautela lo que un \textit{ranking} dice y lo que no dice. Con acumulación real, Pearson obtiene el mejor resultado; con degradación paramétrica, lo hacen conjuntamente coseno, euclídea, HQI y SAM. Pero ninguna diferencia entre métricas resulta significativa tras la corrección de Holm, así que ese orden debe leerse como tendencia, no como una superioridad demostrada. Hacer análisis de datos no es solo producir métricas, \textit{rankings} o \textit{benchmarks}: es también saber hasta dónde sostienen una conclusión. Me llevo de aquí el hábito de no quedarme con un único indicador y de mirar en conjunto el tamaño de la muestra, la cobertura de los datos, el diseño experimental y la significancia estadística antes de afirmar que una métrica es mejor que otra.

A nivel profesional, el TFM me ha ayudado a precisar el tipo de problemas en los que quiero seguir formándome. Disfruté especialmente las partes de construcción del \textit{pipeline}, transformación de datos, automatización, validación, reproducibilidad y análisis de resultados. Al mismo tiempo, el proyecto me confirmó que no me basta con construir sistemas que muevan o procesen datos: necesito entender el problema que hay detrás y tener una forma objetiva de comprobar si la solución funciona. Me sitúo, por tanto, en la intersección entre ingeniería de datos y análisis aplicado, combinando la construcción de infraestructura de procesamiento con la interpretación crítica de lo que esa infraestructura produce.

Si tuviera que resumir el aprendizaje del TFM en una sola idea, sería esta: el valor de una solución de Big Data no depende de la complejidad de su arquitectura, sino de si esa arquitectura, esos datos, esas métricas y ese diseño experimental son realmente los adecuados para responder a la pregunta planteada. El proyecto arrancó con una propuesta tecnológicamente ambiciosa y terminó mostrando que, a esta escala, pandas superaba a Spark, y que un \textit{pipeline} técnicamente impecable no basta si los datos de referencia no permiten evaluarlo. Aprender a construir el \textit{pipeline} fue necesario; aprender a cuestionar si el \textit{pipeline}, los datos y el diseño eran los correctos fue, con diferencia, lo más valioso.
